\section{Lista 3 (04/05/2026)}
%Vou listar os que acho legais

\subsection{Convergência de variáveis aleatórias}
%Duas formas de resolver esse com sparsification e com lemma de Kroenecher e teorema das duas séries.
\begin{exercise}
    Seja $(X_n)_{n\geq1}$ uma sequência de r.v's e $C > 0$ uma constante tal que 
    $\Ex X_n^2 < C$ para todo $n\geq1$ e $\text{Cov}(X_n, X_m) = 0$ para todo $n \neq m$. Prove que 
    \begin{enumerate}[label=(\alph*)]
        \item 
        \[\lim_{n\to\infty} \frac{S_{n^2} - \Ex S_{n^2}}{n^2} = 0 \quad \text{a.s}.\]
        \item 
        \[\lim_{n\to\infty} \frac{S_{n} - \Ex S_{n}}{n} = 0 \quad \text{a.s}.\]
    \end{enumerate}
\end{exercise}

%Esse aqui é bem maneiro
\begin{exercise}
    Sejam $(X_n)_n$ r.v's tais que cada $X_n$ toma valores em um conjunto finito $A_n \subset \R$. Suponha também 
    que $\lim_{n\to\infty} \max_{x \in A_n} \Pr(X_n = x) = 0$. Mostre que
    \[
        F_{X_n}(X_n) \to U[0,1]
    \]
    em distribuição.
\end{exercise}

%Os outros valem a pena fazer para treinar bem as definições.

\subsection{Lei dos Grandes Números}
%Weierstrass já fiz e refiz
%2 e 3 parecem divertidos

\subsection{Teorema Central do Limite}
%novamente 2 e 3 parecem legais

\pagebreak